% Глава: 22-docker-compose
% Статус: наполнено (первый прогон).
\chapter{Docker Compose}
\label{ch:22-docker-compose}

\section{Зачем Compose}

Когда приложение --- это несколько контейнеров
(веб, база данных, очередь), описывать их
отдельными \cmd{docker run} неудобно.
\textbf{Docker Compose} описывает всю систему в одном
YAML-файле и запускает её одной командой.

Оркеструем такие системы на своей ВМ: Compose поднимает
все контейнеры на \fpath{pglab.nsu.ru} одной командой.

\section{Анатомия \fpath{docker-compose.yml}}

\begin{minted}[bgcolor=gray!6]{bash}
services:
  web:
    build: .
    ports:
      - "8080:80"
    depends_on:
      - db
    environment:
      - DATABASE_URL=postgres://db:5432
  db:
    image: postgres:16
    volumes:
      - dbdata:/var/lib/postgresql/data
    environment:
      POSTGRES_PASSWORD: secret

volumes:
  dbdata:
\end{minted}

Ключевые поля:
\begin{center}
\begin{tabular}{lp{0.62\linewidth}}
\toprule
\textbf{Поле} & \textbf{Назначение} \\
\midrule
\cmd{services}      & словарь сервисов (контейнеров) \\
\cmd{image}         & готовый образ \\
\cmd{build}         & собрать из Dockerfile (путь) \\
\cmd{ports}         & проброс портов (host:container) \\
\cmd{volumes}       & тома (именованные или bind mount) \\
\cmd{environment}   & переменные окружения \\
\cmd{depends\_on}   & порядок запуска (запустит сначала) \\
\cmd{networks}      & пользовательские сети \\
\cmd{command}       & переопределить CMD образа \\
\bottomrule
\end{tabular}
\end{center}

\section{Основные команды}

\begin{minted}[bgcolor=gray!6]{bash}
docker compose up            # собрать при необходимости и запустить все сервисы
docker compose up -d         # запустить в фоне
docker compose ps            # список работающих сервисов
docker compose logs web      # логи сервиса "web"
docker compose exec web bash # открыть оболочку в сервисе "web"
docker compose stop          # остановить все
docker compose down          # остановить и удалить контейнеры
docker compose down -v       # также удалить тома, осторожно!
\end{minted}

\section{Сети и DNS}

Компоненты в одном compose-файле автоматически
попадают в одну сеть; имя сервиса --- его DNS-имя:
веб обращается к базе как \cmd{db} (не \cmd{localhost}).

\section{Типичные ошибки}
\begin{itemize}
  \item \cmd{depends\_on} --- только про
    \emph{запуск}, не про \emph{готовность}: база может
    ещё не быть готова принимать подключения, когда
    стартует веб; используйте \cmd{healthcheck} или
    retry в коде;
  \item Порт занят на хосте --- \cmd{docker compose ps}
    покажет; смените host-порт в \cmd{ports};
  \item \cmd{docker compose down -v} удаляет тома и
    данные --- используйте с осторожностью
    (\nameref{ch:21-docker-data});
  \item Несовместимые версии \cmd{compose} и
    \cmd{docker} --- обновите инструмент.
\end{itemize}

\section{Связь с остальным курсом}

Compose комбинирует Dockerfile
(\nameref{ch:20-dockerfile}) и тома
(\nameref{ch:21-docker-data}); ограничения ресурсов
для каждого сервиса --- \nameref{ch:23-docker-resources};
запуск нескольких контейнеров на хосте ---
\nameref{ch:19-docker-basics}.

\section{Задачи для самопроверки}
\begin{enumerate}
  \item Чем Compose удобнее, чем несколько
    \cmd{docker run}?
  \item Какие ключевые поля у сервиса в
    \fpath{docker-compose.yml}?
  \item Как запустить систему в фоне и посмотреть
    логи одного сервиса?
  \item Чем \cmd{depends\_on} не гарантирует
    (и что нужно добавить)?
  \item Как сервисам обращаться друг к другу по
    имени?
\end{enumerate}
